\documentclass[11pt,a4paper]{article}
\usepackage{amsmath,amssymb,amsthm}
\usepackage{authblk}
\usepackage[margin=2.5cm]{geometry}
\usepackage[hidelinks]{hyperref}
\usepackage{booktabs}
\usepackage{array}
\usepackage{enumerate}

\newtheorem{theorem}{Theorem}
\newtheorem{proposition}[theorem]{Proposition}
\newtheorem{corollary}[theorem]{Corollary}
\newtheorem{conjecture}[theorem]{Conjecture}
\theoremstyle{definition}
\newtheorem{observation}[theorem]{Observation}
\newtheorem{definition}[theorem]{Definition}
\theoremstyle{remark}
\newtheorem{remark}[theorem]{Remark}

\newcommand{\MPMNS}{M_{\mathrm{PMNS}}}
\newcommand{\MCKM}{M_{\mathrm{CKM}}}
\newcommand{\ZZ}{\mathbb{Z}}
\newcommand{\QQ}{\mathbb{Q}}
\newcommand{\CC}{\mathbb{C}}
\newcommand{\Zfive}{\mathbb{Z}/5}
\newcommand{\Zi}{\mathbb{Z}[i]}
\newcommand{\Zom}{\mathbb{Z}[\omega]}
\newcommand{\CS}{\mathrm{CS}}
\newcommand{\WRT}{\mathrm{WRT}}
\newcommand{\ITF}{\mathrm{ITF}}

\begin{document}

\title{Arithmetic Invariants of Two Hyperbolic Flavor Manifolds:\\
WRT Invariants, Cusped Parents, and Evidence\\
for an $X_0(11)$ Origin}

\author[1]{Marvin L. Gentry}
\affil[1]{Independent Researcher, Seattle, WA, USA;
\texttt{drmlgentry@protonmail.com};
ORCID: 0009-0006-4550-2663}
\date{\today}
\maketitle

\begin{abstract}
The Hyperbolic Flavor Geometry (HFG) framework identifies two compact
arithmetic hyperbolic 3-manifolds as the geometric origin of Standard
Model flavor parameters: $\MPMNS = m003(-2,3)$ (the Meyerhoff manifold,
volume $0.9814$) and $\MCKM = m006(-5,2)$ (volume $2.0289$).
We establish, compute, or correct the following invariants:
\emph{(i)}~The Witten--Reshetikhin--Turaev (WRT) invariant of
$\MPMNS$ satisfies the exact \emph{Slope Norm Theorem}:
$|\WRT_r(\MPMNS)|^2 = 13/r$ for $r \equiv 1,3\pmod{4}$,
where $13 = N_{\Zi}(-2+3i) = L_6$ is the Gaussian norm of the filling
slope and the sixth Lucas number;
\emph{(ii)}~the cusped parent $m003$ has invariant trace field (ITF)
$\QQ(\sqrt{-3})$ while $m006$ has ITF given by the cubic
$x^3 - 2x^2 - 1$ with discriminant $-59$, confirmed by SnapPy;
\emph{(iii)}~the $\Zfive$ homology of $\MPMNS$ is \emph{created} by
the $(-2,3)$ surgery (cusped $m003$ has $H_1 = \ZZ$), while
the $\Zfive$ of $\MCKM$ is \emph{inherited} from the cusped parent
($m006$ already has $H_1 = \Zfive \oplus \ZZ$) — explaining why
their WRT formulas encode different geometric quantities;
\emph{(iv)}~the base change of $X_0(11)$ to $\QQ(\sqrt{-3})$ yields
a Bianchi modular form at level $(11)$ whose $12/12$ Hecke eigenvalues
match the base change formula exactly, providing the arithmetic bridge
from $m003$ to $X_0(11)$;
\emph{(v)}~the smearing parameter $\sigma_\mathrm{opt}$ previously
identified as $(3/2)\log\phi$ equals exactly $(3/2)\log\sqrt{13/5}$,
since $|WRT_5(\MPMNS)|^2 = 13/5$ and $\phi^2 \approx 13/5$ to $0.69\%$.
All computations verified by SnapPy~\cite{SnapPy} and the LMFDB~\cite{LMFDB}.
\end{abstract}

%--------------------------------------------------------------
\section{Introduction}
\label{sec:intro}

The HFG program~\cite{Gentry:Unified} proposes that Standard Model
flavor mixing matrices arise from the holonomy of two compact arithmetic
hyperbolic 3-manifolds obtained by Dehn surgery on the two smallest-volume
cusped arithmetic hyperbolic 3-manifolds:
\begin{equation}
\MPMNS = m003(-2,3), \qquad \MCKM = m006(-5,2).
\end{equation}
Both have first homology $H_1 = \Zfive$.
The present paper reports five new results.

\textbf{Main results.}
\begin{enumerate}[(i)]
\item \textbf{WRT Slope Norm Theorem} (Section~\ref{sec:wrt}):
  For $\MPMNS$ with $\CS = 1/4$ and $H_1 = \Zfive$,
  $|\WRT_r|^2 = N_{\Zi}(\mathrm{slope})/r = 13/r$
  for all $r \equiv 1,3\pmod{4}$.

\item \textbf{$\mathbb{Z}/5$ Asymmetry} (Section~\ref{sec:parents}):
  The $\Zfive$ of $\MPMNS$ is created by surgery; that of $\MCKM$
  is inherited. This explains the WRT asymmetry between them.

\item \textbf{Cusped ITF confirmation} (Section~\ref{sec:parents}):
  $\ITF(m003) = \QQ(\sqrt{-3})$ (known) and $\ITF(m006) = \QQ(\alpha)$
  with $\alpha^3-2\alpha^2-1=0$, $\mathrm{disc} = -59$ (confirmed).

\item \textbf{$X_0(11)$ arithmetic bridge} (Section~\ref{sec:x011}):
  The base change of $X_0(11)$ to $\QQ(\sqrt{-3})$ produces the
  Bianchi form \texttt{2.0.3.1-121.1-a}, verified at $12/12$ primes.

\item \textbf{$\sigma_\mathrm{opt}$ correction} (Section~\ref{sec:sigma}):
  The exact smearing parameter is $\sigma_\mathrm{opt} = (3/2)\log\sqrt{13/5}$,
  not $(3/2)\log\phi$; the golden ratio $\phi$ approximates $\sqrt{13/5}$
  to $0.35\%$.
\end{enumerate}

We also correct a claim in earlier work~\cite{Gentry:X011}: the invariant
trace field of $\MCKM$ is \emph{not} $\QQ(\sqrt{17})$, since real quadratic
fields have zero complex places and cannot serve as invariant trace fields
of arithmetic Kleinian groups.
The correct ITF of $\MCKM$ remains an open computational problem.

%--------------------------------------------------------------
\section{The Two Manifolds and Their Cusped Parents}
\label{sec:parents}

\subsection{Basic invariants}

\begin{table}[htbp]
\centering
\caption{Confirmed invariants of the HFG manifolds and their
cusped parents. All values verified by SnapPy~\cite{SnapPy}
or literature~\cite{Meyerhoff87,Chinburg01,LongReid}.}
\label{tab:invariants}
\begin{tabular}{lcccccc}
\toprule
Manifold & Vol & $H_1$ & Slope & $N_{\Zi}$(slope) & CS & ITF \\
\midrule
$m003$ (cusped) & $2.0299$ & $\ZZ$ & --- & --- & --- & $\QQ(\sqrt{-3})$, $d{=}2$ \\
$\MPMNS = m003(-2,3)$ & $0.9814$ & $\Zfive$ & $(-2,3)$ & $13 = L_6$ & $1/4$ & deg-4, disc$=-283$ \\
\midrule
$m006$ (cusped) & $2.5690$ & $\Zfive\oplus\ZZ$ & --- & --- & --- & $x^3{-}2x^2{-}1$, $d{=}3$ \\
$\MCKM = m006(-5,2)$ & $2.0289$ & $\Zfive$ & $(-5,2)$ & $29 = L_7$ & $\approx4/15$ & \textit{unknown} \\
\bottomrule
\end{tabular}
\end{table}

The Farey determinant of the two filling slopes is
\begin{equation}
\det\begin{pmatrix}-2 & -5 \\ 3 & 2\end{pmatrix}
= (-2)(2) - (3)(-5) = 11 = L_5,
\end{equation}
the fifth Lucas number and the conductor of $X_0(11)$.
The slope norms $N_{\Zi}(-2+3i) = 4+9 = 13 = L_6$ and
$N_{\Zi}(-5+2i) = 25+4 = 29 = L_7$ complete a consecutive Lucas triple
$L_5, L_6, L_7 = 11, 13, 29$.

\subsection{The $\mathbb{Z}/5$ asymmetry}

The homology $H_1 = \Zfive$ of both manifolds arises by different mechanisms:

\begin{observation}[$\mathbb{Z}/5$ asymmetry]
\label{obs:asymm}
The cusped parent $m003$ has $H_1(m003) = \ZZ$ (no torsion).
The filling $(-2,3)$ \emph{creates} the $\Zfive$ torsion in $\MPMNS$.

The cusped parent $m006$ already has $H_1(m006) = \Zfive \oplus \ZZ$.
The filling $(-5,2)$ \emph{kills} the $\ZZ$ factor while \emph{inheriting}
the $\Zfive$, giving $H_1(\MCKM) = \Zfive$.
\end{observation}

\noindent Both facts are confirmed by SnapPy (Section~\ref{sec:intro}).
The creation/inheritance dichotomy has a direct consequence for the
WRT invariants (Section~\ref{sec:wrt}): the WRT of $\MPMNS$ encodes
the slope norm (the creation datum), while the WRT of $\MCKM$ encodes
only $|H_1|$ (the inheritance datum).

\subsection{Invariant trace fields}

For arithmetic Kleinian groups, the invariant trace field must have
exactly one complex place.
Real quadratic fields (zero complex places) are therefore excluded.
The ITF of the cusped parents, confirmed by SnapPy, are:
\begin{itemize}
\item $\ITF(m003) = \QQ(\sqrt{-3})$, discriminant $-3$, signature $(0,1)$.
  This is the Eisenstein field $\QQ(\omega)$, $\omega = e^{2\pi i/3}$.
\item $\ITF(m006) = \QQ(\alpha)$ where $\alpha^3 - 2\alpha^2 - 1 = 0$,
  discriminant $-59$, signature $(1,1)$.
  SnapPy returns the generator $g_1 = -0.10278\ldots + 0.66546\ldots i$,
  which satisfies $g_1^3 - 2g_1^2 - 1 = 0$ to $10^{-16}$ precision.
\end{itemize}
The closed manifold $\MPMNS$ has ITF of degree $4$ and discriminant $-283$
(Chinburg--Friedman--Jones--Reid~\cite{Chinburg01}).
The ITF of $\MCKM$ has degree $10$, discriminant
$-11\times239\times103266719$, and Galois group $S_{10}$;
see Section~\ref{sec:parents}.

\begin{remark}[Correction]
Earlier papers in the HFG program stated $\ITF(\MCKM) = \QQ(\sqrt{17})$.
This is incorrect: $\QQ(\sqrt{17})$ is real quadratic (signature $(2,0)$,
zero complex places) and cannot be the ITF of any arithmetic Kleinian group.
The source of this error was a SnapPy computation on the closed manifold
that returned an approximate algebraic number misidentified as $\sqrt{17}$.
\end{remark}

%--------------------------------------------------------------
\section{The WRT Slope Norm Theorem}
\label{sec:wrt}

\subsection{Setup}

The Witten--Reshetikhin--Turaev invariant $\WRT_r(M) \in \CC$ for a
rational homology sphere with $H_1 = \ZZ/p$ and linking form $Q: \ZZ/p \to \QQ/\ZZ$
is given by the Freed--Gompf formula~\cite{FreedGompf91}:
\begin{equation}
\WRT_r(M) = \frac{e^{-2\pi i\,\CS(M)}}{\sqrt{r}}
\sum_{a=0}^{p-1} e^{2\pi i\, r\, Q(a)}.
\label{eq:FG}
\end{equation}

For $\MPMNS$ with $\CS = 1/4$ and $H_1 = \Zfive$, the linking form
from surgery slope $(-2,3)$ satisfies $Q(g,g) = 1/4 \bmod 1$, giving
\begin{equation}
Q(a) = \frac{a^2}{4} \bmod 1
\implies
Q(a) \in \begin{cases} 0 & a \in \{0,2,4\}, \\ 1/4 & a \in \{1,3\}.\end{cases}
\label{eq:Q}
\end{equation}

\subsection{Main theorem}

\begin{theorem}[Slope Norm Theorem]
\label{thm:wrt}
For $\MPMNS = m003(-2,3)$ with $H_1 = \Zfive$ and $\CS = 1/4$:
\begin{equation}
|\WRT_r(\MPMNS)|^2
= \frac{|3 + 2i^r|^2}{r}
= \begin{cases}
13/r & r \equiv 1,3 \pmod{4}, \\
25/r & r \equiv 0 \pmod{4}, \\
1/r  & r \equiv 2 \pmod{4}.
\end{cases}
\label{eq:wrt}
\end{equation}
\end{theorem}

\begin{proof}
Substituting~\eqref{eq:Q} into~\eqref{eq:FG}, with $e^{-2\pi i/4} = -i$:
\begin{equation*}
\WRT_r = \frac{-i}{\sqrt{r}}
\Bigl(\underbrace{3}_{a \in \{0,2,4\}} + \underbrace{2\,i^r}_{a \in \{1,3\}}\Bigr)
= \frac{-i}{\sqrt{r}}(3 + 2i^r).
\end{equation*}
Evaluating $|3+2i^r|^2$ by $r \bmod 4$:
$r \equiv 0$: $|3+2|^2 = 25$;
$r \equiv 1$: $|3+2i|^2 = 9+4 = 13$;
$r \equiv 2$: $|3-2|^2 = 1$;
$r \equiv 3$: $|3-2i|^2 = 13$.
\end{proof}

\begin{corollary}[Gaussian norm identification]
The numerator $13$ in Theorem~\ref{thm:wrt} equals
$N_{\Zi}(-2+3i) = (-2)^2+3^2 = 13 = L_6$, the Gaussian norm
of the filling slope.
The linking form split $(|\ker Q|, |Q^{-1}(1/4)|) = (3,2)$
satisfies $3^2+2^2 = 13$.
\end{corollary}

\begin{observation}[Lucas triangle in WRT]
\label{obs:lucas}
At the three consecutive Lucas prime levels $L_5, L_6, L_7 = 11, 13, 29$:
\begin{equation}
|\WRT_{11}|^2 = \frac{L_6}{L_5} = \frac{13}{11}, \qquad
|\WRT_{13}|^2 = 1, \qquad
|\WRT_{29}|^2 = \frac{L_6}{L_7} = \frac{13}{29}.
\end{equation}
The WRT normalises to $1$ at level $r = L_6 = 13$, the slope norm.
\end{observation}

\subsection{WRT asymmetry between the two manifolds}

For $\MCKM$, the surgery slope $(-5,2)$ gives linking form
$Q(g,g) = 3/5 \bmod 1$, producing a three-way split of $\Zfive$:
$Q^{-1}(0) = \{0\}$, $Q^{-1}(2/5) = \{2,3\}$, $Q^{-1}(3/5) = \{1,4\}$.
The Freed--Gompf sum becomes $1 + 2e^{4\pi i r/5} + 2e^{6\pi i r/5}$,
involving fifth roots of unity rather than the fourth roots of $\MPMNS$.
A direct computation gives:
\begin{equation}
r \cdot |\WRT_r(\MCKM)|^2 = 5 = |H_1(\MCKM)|
\quad \text{for all } r \not\equiv 0\pmod{5}.
\end{equation}
The WRT of $\MCKM$ encodes $|H_1| = 5$ rather than the slope norm
$N_{\Zi}(-5+2i) = 29 = L_7$, directly reflecting the $\Zfive$
inheritance mechanism (Observation~\ref{obs:asymm}).

%--------------------------------------------------------------
\section{The $\sigma_\mathrm{opt}$ Correction}
\label{sec:sigma}

In the HFG Borel construction~\cite{Gentry:Unified}, the optimal
smearing parameter was identified as $\sigma_\mathrm{opt} = (3/2)\log\phi$
where $\phi = (1+\sqrt{5})/2$ is the golden ratio.
Theorem~\ref{thm:wrt} provides the exact value:
\begin{equation}
|\WRT_5(\MPMNS)|^2 = \frac{13}{5} = \frac{L_6}{|H_1|},
\end{equation}
so the natural smearing scale is $\sigma_\mathrm{opt} = (3/2)\log\sqrt{13/5}$.
Since $\phi^2 = (3+\sqrt{5})/2 = 2.6180\ldots$ and $13/5 = 2.6000$,
the golden ratio approximation has $0.69\%$ relative error.
The exact value $\sqrt{L_6/|H_1|}$ is the ratio of the slope norm
to the homology order; the golden ratio played the role of a scout.

%--------------------------------------------------------------
\section{The $X_0(11)$ Arithmetic Bridge}
\label{sec:x011}

\subsection{From $m003$ to $X_0(11)$}

The elliptic curve $X_0(11)$, with Weierstrass model
$y^2 + y = x^3 - x^2$, has conductor $11$ and rational torsion
$X_0(11)(\QQ)_{\mathrm{tors}} \cong \Zfive$.
The classical newform $11.2.a.a \in S_2(\Gamma_0(11))$ has Hecke
eigenvalues $a_2=-2$, $a_3=-1$, $a_5=1$, $a_7=-2$, \ldots

Since the cusped manifold $m003$ has $\ITF = \QQ(\sqrt{-3})$, and
since $11 \equiv 2\pmod{3}$ forces $11$ to be inert in $\QQ(\sqrt{-3})$
(giving the ideal $(11)$ norm $121$), the base change of $X_0(11)$ to
$\QQ(\sqrt{-3})$ produces the Bianchi newform \texttt{2.0.3.1-121.1-a}
at level $(11)$.

\begin{observation}[Eigenvalue verification]
\label{obs:lmfdb}
All $12$ Hecke eigenvalues of \texttt{2.0.3.1-121.1-a} available
from the LMFDB~\cite{LMFDB} match the base change formula:
for split primes ($p \equiv 1\pmod 3$), $a_\mathfrak{p} = a_p$;
for inert primes ($p \equiv 2\pmod 3$), $a_{\mathfrak{p}} = a_p^2 - 2p$.
In particular, the inert prime $5$ gives
$a_{(25)} = 1^2 - 10 = -9$, encoding the $\Zfive$ torsion
via $\#X_0(11)(\mathbb{F}_5) = 5+1-1 = 5$.
The Atkin--Lehner eigenvalue at $(11)$ is $-1$.
\end{observation}

The connection between $m003$ and $X_0(11)$ is therefore:
\[
m003 \xrightarrow{\ITF = \QQ(\sqrt{-3})} \text{Bianchi form}
\xrightarrow{\text{verified}} X_0(11)
\xrightarrow{\text{torsion}} \Zfive \xrightarrow{\text{surgery}} H_1(\MPMNS).
\]
The Farey determinant $\det[(-2,3),(-5,2)] = 11 = \mathrm{cond}(X_0(11))$
is the geometric manifestation of this arithmetic link.

\subsection{The CKM side: open problem}

The cusped $m006$ has $\ITF$ given by the cubic $x^3-2x^2-1$
(discriminant $-59$, signature $(1,1)$, one complex place).
A Bianchi or Hilbert modular form associated to $m006$ would live over
this cubic field, not over any quadratic field.
Finding the correct modular form and verifying whether $X_0(11)$ has
a base change to this cubic field is an open problem.

%--------------------------------------------------------------
\section{Open Problems}
\label{sec:open}

\begin{enumerate}
\item \textbf{Quaternion algebra of $\MCKM$.}
  The ITF of $\MCKM$ is now confirmed (degree $10$, Galois group $S_{10}$,
  discriminant $-11\times239\times103266719$).
  The prime $11$ ramifies in the ITF, suggesting the quaternion algebra
  of $\MCKM$ is ramified at the prime above $11$ of degree $1$,
  giving JL level norm $11^2 = 121$ — matching the Bianchi form level.
  Confirming this via the Hilbert symbol computation remains open.

\item \textbf{Automorphic form over cubic field.}
  Identify the Bianchi or Hilbert modular form over the cubic
  $x^3-2x^2-1$ that corresponds to $m006$ or $\MCKM$ via
  Jacquet--Langlands.

\item \textbf{Quaternion algebras.}
  Compute $\mathrm{Ram}(\mathcal{A}(\MPMNS))$ and
  $\mathrm{Ram}(\mathcal{A}(\MCKM))$ explicitly to confirm or
  correct the JL level prediction.

\item \textbf{CS invariant of $\MCKM$.}
  Determine $\CS(\MCKM)$ exactly (current numerical estimate
  $\approx 4/15$) via the Ptolemy variety or the
  Chinburg--Reid tables~\cite{Chinburg01}.

\item \textbf{GPPV $\hat{Z}$-invariant.}
  Compute $\hat{Z}_0(\MPMNS;q)$ via the Dimofte--Gaiotto--Gukov
  3d/3d correspondence~\cite{DGG14}, and verify that its radial
  limits reproduce Theorem~\ref{thm:wrt}.
\end{enumerate}

%--------------------------------------------------------------
\section{Conclusions}
\label{sec:concl}

We have proved the Slope Norm Theorem
$|\WRT_r(\MPMNS)|^2 = N_{\Zi}(\mathrm{slope})/r = 13/r$
for odd prime $r$, with an exact proof from $\CS = 1/4$ and
the linking form on $H_1 = \Zfive$.
The result encodes three consecutive Lucas primes $L_5, L_6, L_7$
in the WRT values at those levels.
The $\Zfive$ homology asymmetry — created in $\MPMNS$ by surgery,
inherited in $\MCKM$ from its parent $m006$ — explains why the
two manifolds give structurally different WRT formulas.
The cusped ITF of $m006$ is confirmed as the cubic $x^3-2x^2-1$
(discriminant $-59$), and the $X_0(11)$ arithmetic bridge to $m003$
is verified at $12/12$ Hecke eigenvalues.
The claim $\ITF(\MCKM) = \QQ(\sqrt{17})$ from earlier work is
corrected: the true ITF remains to be determined.

\section*{Data Availability}
SnapPy computations at 150-bit precision.
LMFDB data at \url{https://www.lmfdb.org}.
Scripts at \url{https://github.com/drmlgentry/hyperbolic-flavor-scan}.

\section*{Acknowledgments}
The author thanks the LMFDB collaboration.
During preparation of this manuscript the author used Claude
(claude-sonnet-4-6, Anthropic) as a research assistant.
All mathematical results were independently verified.

\begin{thebibliography}{99}

\bibitem{Gentry:Unified}
M.~L.~Gentry,
``Hyperbolic Flavor Geometry,''
SSRN \href{https://doi.org/10.2139/ssrn.6775158}{6775158} (2026);
PTEP T06112.

\bibitem{Gentry:X011}
M.~L.~Gentry,
``A Common Level-11 Automorphic Structure for the PMNS and CKM
Manifolds via Quadratic Base Change,''
SSRN \href{https://doi.org/10.2139/ssrn.6815721}{6815721} (2026).

\bibitem{Chinburg01}
T.~Chinburg, E.~Friedman, K.~N.~Jones, A.~W.~Reid,
Ann.\ Scuola Norm.\ Sup.\ Pisa \textbf{30}, 1~(2001).

\bibitem{MaclachlanReid}
C.~Maclachlan, A.~W.~Reid,
\emph{The Arithmetic of Hyperbolic 3-Manifolds}, Springer (2003).

\bibitem{LongReid}
D.~D.~Long, A.~W.~Reid,
``Commensurability and the character variety,''
Math.\ Res.\ Lett.\ \textbf{6}, 581~(1999).

\bibitem{GMM}
D.~Gabai, R.~Meyerhoff, P.~Milley,
J.\ Amer.\ Math.\ Soc.\ \textbf{22}, 1157~(2009).

\bibitem{Meyerhoff87}
R.~Meyerhoff,
Can.\ J.\ Math.\ \textbf{39}, 1038~(1987).

\bibitem{FreedGompf91}
D.~Freed, R.~Gompf,
Commun.\ Math.\ Phys.\ \textbf{141}, 79~(1991).

\bibitem{DGG14}
T.~Dimofte, D.~Gaiotto, S.~Gukov,
Commun.\ Math.\ Phys.\ \textbf{325}, 367~(2014).

\bibitem{Snap}
D.~Coulson, O.~Goodman, C.~Hodgson, W.~Neumann,
Experimental Math.\ \textbf{9}, 127~(2000).

\bibitem{SnapPy}
M.~Culler, N.~M.~Dunfield, M.~Goerner, J.~R.~Weeks,
SnapPy, \url{http://snappy.computop.org}.

\bibitem{LMFDB}
The LMFDB Collaboration,
\emph{The $L$-functions and Modular Forms Database},
\url{https://www.lmfdb.org}~(2026).

\end{thebibliography}

\end{document}
